\uuid{62Ow}
\titre{ Implémentation de ReLU et de sa dérivée }
\niveau{L3}
\module{Analyse de données}
\chapitre{Réseaux de neurones}
\sousChapitre{Fonctions d'activation, rétropropagation}
\theme{ReLU, dérivée de ReLU, neurone mort, backward pass}
\auteur{Maxime Nguyen}
\datecreate{2026-03-19}
\organisation{AMSCC}
\difficulte{2}

\contenu{
\texte{
  On compare trois implémentations de ReLU et de sa dérivée utilisée dans le backward pass :
  \begin{itemize}
    \item[A.] \texttt{relu\_backward} met à zéro \texttt{dZ} là où \texttt{Z <= 0}
    \item[B.] \texttt{relu\_backward} retourne \texttt{dA * Z}
    \item[C.] \texttt{relu\_backward} met à zéro \texttt{dZ} là où \texttt{Z < 0}
  \end{itemize}
}
\begin{enumerate}

\item
  \question{Laquelle est correcte ?}
  \reponse{
    \textbf{A} est la référence standard : $\text{ReLU}'(z) = \mathbf{1}_{z > 0}$,
    donc on annule le gradient là où $z \leq 0$, ce qui inclut $z = 0$.
  }

\item
  \question{Quelle est la différence entre A et C ? Est-ce important en pratique ?}
  \reponse{
    A met à zéro pour $z \leq 0$ (inclut $z = 0$), C met à zéro pour $z < 0$ (exclut $z = 0$).
    En $z = 0$, la dérivée de ReLU n'est mathématiquement pas définie (point de non-dérivabilité).
    En pratique, la différence est négligeable : la probabilité qu'une activation soit exactement $0$
    est quasi nulle sur des données réelles, et les deux conventions donnent des résultats identiques.
  }

\item
  \question{Qu'est-ce qui ne va pas dans B ?}
  \reponse{
    B retourne \texttt{dA * Z}, c'est-à-dire le produit du gradient amont par la valeur de la pré-activation $z$,
    au lieu de le masquer par $\mathbf{1}_{z > 0}$.
    Cette formule est incorrecte : pour $z > 0$, on devrait passer le gradient tel quel (\texttt{dA}),
    et non le multiplier par $z$. B introduit donc une erreur dans le calcul du gradient.
  }

\end{enumerate}
}
